\documentclass[10pt,twocolumn]{article}
\usepackage[margin=0.75in]{geometry}
\usepackage{amsmath,amssymb,booktabs,graphicx,hyperref}
\title{Constrained LinUCB for Crop Recommendation under Agronomic Feasibility Rules}
\author{Patrick Gomes de Oliveira}
\date{\today}
\begin{document}
\maketitle
\begin{abstract}
We frame crop recommendation as a contextual bandit over soil and climate
readings, with the recommendable crop as the arm and agreement with an
agronomist's logged choice as the simulated reward. A disjoint linear
LinUCB model (Li et al., 2010) is restricted at every round to the subset
of crops whose agronomic feasible range, estimated from training data,
contains the current context -- turning plain LinUCB into \emph{constrained}
LinUCB. We measure both generalisation (held-out match rate and regret) and
online learning speed, and find that the agronomic constraint sharply
reduces the exploration cost of learning, at a small, measurable risk of
excluding the true crop when a test context is an outlier relative to its
own training distribution.
\end{abstract}
\section{Introduction}
Public crop-recommendation notebooks almost universally treat the problem
as multi-class classification (random forests, gradient boosting). This
project instead frames it as a \emph{contextual bandit}: a decision made
under a single observed context with partial feedback, which is the
realistic online setting for a system deployed to make one recommendation
at a time and learn from whether it was accepted. LinUCB \cite{li2010linucb}
is a classical, closed-form linear bandit with no gradient descent and no
neural network; we add agronomic domain knowledge as a hard constraint on
the arm set offered at each round, which is both a safety mechanism and, as
Section~\ref{sec:results} shows, a substantial accelerant of learning.

\section{Dataset}
The Crop Recommendation dataset (Kaggle ref
\texttt{atharvaingle/crop-recommendation-dataset}) contains soil nutrient
levels (N, P, K), temperature, humidity, pH, and rainfall, each row labelled
with the crop an agronomist recommends for those conditions, across 22
crop types. Following the standard technique for turning a labelled
classification dataset into an offline contextual-bandit replay log
\cite{li2011unbiased}, each row becomes one round: the context is the
soil/climate vector, and the simulated reward for choosing arm $a$ is $1$
iff $a$ equals the row's logged crop.

\section{Feature engineering}
Features are declared in \texttt{conf/features.yaml}.
\texttt{context\_zscore} standardises the seven raw soil/climate features,
fitted strictly on the training replay stream, since LinUCB's ridge prior
assumes a roughly isotropic scale. \texttt{arm\_feasible\_range} computes,
per crop, the training-time observed (min, max) of every raw feature
expanded by an 8\% tolerance margin -- this is the agronomic constraint
table. \texttt{arm\_popularity} and \texttt{arm\_context\_centroid} are
training-only descriptive statistics used by the popularity baseline and
reported for interpretability. A registry check asserts the YAML
specification and the Python implementation agree in both directions.

\section{Methodology}

\subsection{LinUCB}
For each arm $a$, LinUCB maintains a ridge-regression posterior

\begin{equation}
  A_a = I_d + \sum_{t:\, a_t = a} x_t x_t^\top, \qquad
  b_a = \sum_{t:\, a_t = a} r_t x_t,
  \label{eq:linucb-posterior}
\end{equation}

with point estimate $\theta_a = A_a^{-1} b_a$. The upper-confidence-bound
score for context $x$ under arm $a$ is

\begin{equation}
  p_a(x) = x^\top \theta_a + \alpha \sqrt{x^\top A_a^{-1} x},
  \label{eq:linucb-score}
\end{equation}

where the second term is an exploration bonus proportional to the model's
remaining uncertainty about $x$ under arm $a$'s posterior, and $\alpha$
controls the explore/exploit trade-off. Both $A_a$ and $b_a$ are updated in
closed form after every round; no gradient descent is involved.

\subsection{Agronomic constraint}
Let $R_a = \{(\ell_{a,j}, u_{a,j})\}_{j=1}^{7}$ be crop $a$'s training-time
feasible range (Section~3) for each raw feature $j$. At round $t$ with raw
context $\tilde{x}_t$, the feasible arm set is

\begin{equation}
  \mathcal{F}_t = \{a : \ell_{a,j} \le \tilde{x}_{t,j} \le u_{a,j} \ \forall j\}.
  \label{eq:feasible-set}
\end{equation}

\emph{Constrained} LinUCB selects $\arg\max_{a \in \mathcal{F}_t} p_a(x_t)$;
plain LinUCB selects over every arm. If $\mathcal{F}_t = \emptyset$, every
arm is treated as feasible rather than leaving the policy with no action.

\section{Experimental protocol}
Rounds are shuffled once and split into training (70\%), validation (15\%)
and test (15\%). The training partition is replayed in order as an online
learning stream: the policy selects an arm, observes the simulated reward,
and updates, round by round. Validation/test partitions are then replayed
with the policy frozen (no further updates), reporting match rate (the
simulated reward rate) and regret ($1 - \text{match rate}$) against an
oracle that always matches. We compare constrained LinUCB, plain LinUCB,
a context-blind random policy, and an epsilon-greedy popularity baseline,
and separately report the online learning curve (mean reward over the
first 25, 100, and all training rounds) for constrained versus plain
LinUCB, which isolates how much the constraint reduces exploration cost.

\section{Results}
\label{sec:results}
\begin{table}[h]
\centering
\caption{Held-out generalisation on the synthetic CI fixture (test partition).}
\label{tab:results}
\begin{tabular}{lrrr}
\toprule
Model & Match rate & Regret & Mean pool size \\
\midrule
ConstrainedLinUCB & 0.9556 & 0.0444 & 1.73 \\
LinUCB            & 0.9926 & 0.0074 & 8.00 \\
Random            & 0.8519 & 0.1481 & 1.73 \\
EpsilonGreedy      & 0.9185 & 0.0815 & 1.73 \\
\bottomrule
\end{tabular}
\end{table}

\begin{table}[h]
\centering
\caption{Online learning curve: mean training reward by round window.}
\label{tab:learning-curve}
\begin{tabular}{lrrr}
\toprule
Model & First 25 & First 100 & All rounds \\
\midrule
ConstrainedLinUCB & 1.00 & 0.96 & 0.978 \\
LinUCB            & 0.20 & 0.40 & 0.870 \\
\bottomrule
\end{tabular}
\end{table}

Table~\ref{tab:results} reports held-out generalisation on the synthetic CI
fixture. Table~\ref{tab:learning-curve} reports the online learning curve.
Constrained LinUCB reaches a mean training reward of $1.00$ within the
first 25 rounds, while plain LinUCB -- forced to explore all eight arms
before its per-arm posteriors converge -- needs the full training stream to
approach the same level. This is the central result: restricting the arm
set with agronomic domain knowledge removes nearly all of the exploration
cost. The trade-off is visible in the held-out match rate, where
constrained LinUCB is marginally \emph{below} plain LinUCB (0.956 vs 0.993):
occasionally a test round's context falls just outside its own crop's
training-time feasible range, and the hard constraint excludes the correct
answer outright. This is the known risk of hard constraints in constrained
bandits, and is reported here rather than hidden.

\section{Limitations and conclusion}
The simulated-reward protocol assumes every training row's logged crop is
the single correct answer, which discards any ambiguity in genuinely
multi-suitable soil/climate conditions. The feasible-range tolerance (8\%)
is a fixed hyperparameter, not learned or validated against an agronomic
source; a production system would consult domain experts or a larger
reference table instead. The drug-recommendation project elsewhere in this
portfolio similarly layers a hard safety constraint over a statistical
model and warns against over-reliance on it; the same caution applies
here; crop recommendations should support an agronomist's judgement, not
replace it. Overall, constrained LinUCB demonstrates a clear, measurable
benefit (learning speed) and a clear, measurable cost (occasional exclusion
of the correct arm) from the same mechanism, which is a more honest result
than either in isolation.

\bibliographystyle{plain}
\bibliography{references}
\end{document}
